\documentclass{openreader}
\title{Ἰωάννου γʹ}
\date{}
\setmainlanguage{persian}
\setmainfont{Noto Naskh Arabic}
\setotherlanguage{hebrew}
\setotherlanguage{greek}
\newfontfamily\hebrewfont[Script=Hebrew]{SBL BibLit}
\newfontfamily\greekfont[Script=Greek]{SBL BibLit}
\FootnoteStyle{lettered-by-verse}
\OpenReadersBiblesLanguage{greek}
\begin{document}
\ORBselectlanguage{greek}
\maketitle
\raggedbottom 
\fontsize{16pt}{24pt}\selectfont

\ChapterHeader{1}{فصل ۱}

\MainTextVerseMark{1}{1}Ὁ πρεσβύτερος Γαΐῳ\FnGenericTemplate{ا}{\Parse{برایی مذکر مفرد} \Form{Γάϊος} \Gloss{گایوس} } τῷ ἀγαπητῷ, ὃν ἐγὼ ἀγαπῶ ἐν ἀληθείᾳ. 
\MainTextVerseMark{1}{2}Ἀγαπητέ, περὶ πάντων εὔχομαί\FnGenericTemplate{ا}{\Parse{حال میانه اخباری اول‌شخص مفرد} \Form{εὔχομαι} \Gloss{دعا کردن} } σε εὐοδοῦσθαι\FnGenericTemplate{ب}{\Parse{حال مجهول مصدر} \Form{εὐοδόομαι} \Gloss{کامیاب بودن} } καὶ ὑγιαίνειν,\FnGenericTemplate{پ}{\Parse{حال معلوم مصدر} \Form{ὑγιαίνω} \Gloss{تندرست بودن} } καθὼς εὐοδοῦταί\FnGenericTemplate{ت}{\Parse{حال مجهول اخباری سوم‌شخص مفرد} \Form{εὐοδόομαι} \Gloss{کامیاب بودن} } σου ἡ ψυχή. 
\MainTextVerseMark{1}{3}ἐχάρην\FnGenericTemplate{ا}{\Parse{ماضی ساده مجهول اخباری اول‌شخص مفرد} \Form{χαίρω} } γὰρ λίαν\FnGenericTemplate{ب}{\Form{λίαν} \Gloss{زیاد} } ἐρχομένων\FnGenericTemplate{پ}{\Parse{حال میانه وجه وصفی اضافی مذکر جمع} \Form{ἔρχομαι} } ἀδελφῶν καὶ μαρτυρούντων\FnGenericTemplate{ت}{\Parse{حال معلوم وجه وصفی اضافی مذکر جمع} \Form{μαρτυρέω} } σου τῇ ἀληθείᾳ, καθὼς σὺ ἐν ἀληθείᾳ περιπατεῖς. 
\MainTextVerseMark{1}{4}μειζοτέραν\FnGenericTemplate{ا}{\Parse{مفعولی مؤنث مفرد} \Form{μειζότερος} \Gloss{بیشتر} } τούτων οὐκ ἔχω ⸀χαράν, ἵνα ἀκούω\FnGenericTemplate{ب}{\Parse{حال معلوم التزامی اول‌شخص مفرد} \Form{ἀκούω} } τὰ ἐμὰ τέκνα ἐν ⸀τῇ ἀληθείᾳ περιπατοῦντα.\FnGenericTemplate{پ}{\Parse{حال معلوم وجه وصفی مفعولی خنثی جمع} \Form{περιπατέω} } 
\MainTextVerseMark{1}{5}Ἀγαπητέ, πιστὸν ποιεῖς ὃ ἐὰν ἐργάσῃ\FnGenericTemplate{ا}{\Parse{ماضی ساده میانه التزامی دوم‌شخص مفرد} \Form{ἐργάζομαι} } εἰς τοὺς ἀδελφοὺς καὶ ⸀τοῦτο ξένους,\FnGenericTemplate{ب}{\Parse{مفعولی مذکر جمع} \Form{ξένος} \Gloss{بیگانه} } 
\MainTextVerseMark{1}{6}οἳ ἐμαρτύρησάν\FnGenericTemplate{ا}{\Parse{ماضی ساده معلوم اخباری سوم‌شخص جمع} \Form{μαρτυρέω} } σου τῇ ἀγάπῃ ἐνώπιον ἐκκλησίας, οὓς καλῶς ποιήσεις\FnGenericTemplate{ب}{\Parse{آینده معلوم اخباری دوم‌شخص مفرد} \Form{ποιέω} } προπέμψας\FnGenericTemplate{پ}{\Parse{ماضی ساده معلوم وجه وصفی فاعلی مذکر مفرد} \Form{προπέμπω} \Gloss{روانۀ سفر کردن} } ἀξίως\FnGenericTemplate{ت}{\Form{ἀξίως} \Gloss{به طور شایسته} } τοῦ θεοῦ· 
\MainTextVerseMark{1}{7}ὑπὲρ γὰρ τοῦ ὀνόματος ἐξῆλθον\FnGenericTemplate{ا}{\Parse{ماضی ساده معلوم اخباری سوم‌شخص جمع} \Form{ἐξέρχομαι} } μηδὲν λαμβάνοντες\FnGenericTemplate{ب}{\Parse{حال معلوم وجه وصفی فاعلی مذکر جمع} \Form{λαμβάνω} } ἀπὸ τῶν ⸀ἐθνικῶν.\FnGenericTemplate{پ}{\Parse{اضافی مذکر جمع} \Form{ἐθνικός} \Gloss{ملت} } 
\MainTextVerseMark{1}{8}ἡμεῖς οὖν ὀφείλομεν ⸀ὑπολαμβάνειν\FnGenericTemplate{ا}{\Parse{حال معلوم مصدر} \Form{ὑπολαμβάνω} \Gloss{پذیرفتن} } τοὺς τοιούτους, ἵνα συνεργοὶ\FnGenericTemplate{ب}{\Parse{فاعلی مذکر جمع} \Form{συνεργός} \Gloss{همکار} } γινώμεθα\FnGenericTemplate{پ}{\Parse{حال میانه التزامی اول‌شخص جمع} \Form{γίνομαι} } τῇ ἀληθείᾳ. 
\MainTextVerseMark{1}{9}Ἔγραψά\FnGenericTemplate{ا}{\Parse{ماضی ساده معلوم اخباری اول‌شخص مفرد} \Form{γράφω} } ⸀τι τῇ ἐκκλησίᾳ· ἀλλ’ ὁ φιλοπρωτεύων\FnGenericTemplate{ب}{\Parse{حال معلوم وجه وصفی فاعلی مذکر مفرد} \Form{φιλοπρωτεύω} \Gloss{دوست داشتن جایگاه نخست} } αὐτῶν Διοτρέφης\FnGenericTemplate{پ}{\Parse{فاعلی مذکر مفرد} \Form{Διοτρέφης} \Gloss{دیوتْرِفیس} } οὐκ ἐπιδέχεται\FnGenericTemplate{ت}{\Parse{حال میانه اخباری سوم‌شخص مفرد} \Form{ἐπιδέχομαι} \Gloss{پذیرفتن} } ἡμᾶς. 
\MainTextVerseMark{1}{10}διὰ τοῦτο, ἐὰν ἔλθω,\FnGenericTemplate{ا}{\Parse{ماضی ساده معلوم التزامی اول‌شخص مفرد} \Form{ἔρχομαι} } ὑπομνήσω\FnGenericTemplate{ب}{\Parse{آینده معلوم اخباری اول‌شخص مفرد} \Form{ὑπομιμνῄσκω} \Gloss{یادآوری کردن} } αὐτοῦ τὰ ἔργα ἃ ποιεῖ, λόγοις πονηροῖς φλυαρῶν\FnGenericTemplate{پ}{\Parse{حال معلوم وجه وصفی فاعلی مذکر مفرد} \Form{φλυαρέω} \Gloss{بد گفتن} } ἡμᾶς, καὶ μὴ ἀρκούμενος\FnGenericTemplate{ت}{\Parse{حال مجهول وجه وصفی فاعلی مذکر مفرد} \Form{ἀρκέω} \Gloss{راضی شدن} } ἐπὶ τούτοις οὔτε αὐτὸς ἐπιδέχεται\FnGenericTemplate{ث}{\Parse{حال میانه اخباری سوم‌شخص مفرد} \Form{ἐπιδέχομαι} \Gloss{پذیرفتن} } τοὺς ἀδελφοὺς καὶ τοὺς βουλομένους\FnGenericTemplate{ج}{\Parse{حال میانه وجه وصفی مفعولی مذکر جمع} \Form{βούλομαι} } κωλύει\FnGenericTemplate{چ}{\Parse{حال معلوم اخباری سوم‌شخص مفرد} \Form{κωλύω} \Gloss{جلوگیری کردن} } καὶ ἐκ τῆς ἐκκλησίας ἐκβάλλει. 
\MainTextVerseMark{1}{11}Ἀγαπητέ, μὴ μιμοῦ\FnGenericTemplate{ا}{\Parse{حال میانه امری دوم‌شخص مفرد} \Form{μιμέομαι} \Gloss{تقلید کردن} } τὸ κακὸν ἀλλὰ τὸ ἀγαθόν. ὁ ἀγαθοποιῶν\FnGenericTemplate{ب}{\Parse{حال معلوم وجه وصفی فاعلی مذکر مفرد} \Form{ἀγαθοποιέω} \Gloss{خوبی کردن} } ἐκ τοῦ θεοῦ ἐστιν· ὁ κακοποιῶν\FnGenericTemplate{پ}{\Parse{حال معلوم وجه وصفی فاعلی مذکر مفرد} \Form{κακοποιέω} \Gloss{بدی کردن} } οὐχ ἑώρακεν\FnGenericTemplate{ت}{\Parse{کامل معلوم اخباری سوم‌شخص مفرد} \Form{ὁράω} } τὸν θεόν. 
\MainTextVerseMark{1}{12}Δημητρίῳ\FnGenericTemplate{ا}{\Parse{برایی مذکر مفرد} \Form{Δημήτριος} \Gloss{دیمیتریوس} } μεμαρτύρηται\FnGenericTemplate{ب}{\Parse{کامل مجهول اخباری سوم‌شخص مفرد} \Form{μαρτυρέω} } ὑπὸ πάντων καὶ ὑπὸ αὐτῆς τῆς ἀληθείας· καὶ ἡμεῖς δὲ μαρτυροῦμεν, καὶ ⸀οἶδας\FnGenericTemplate{پ}{\Parse{کامل معلوم اخباری دوم‌شخص مفرد} \Form{οἶδα} } ὅτι ἡ μαρτυρία ἡμῶν ἀληθής\FnGenericTemplate{ت}{\Parse{فاعلی مؤنث مفرد} \Form{ἀληθής} \Gloss{حقیقی} } ἐστιν. 
\MainTextVerseMark{1}{13}Πολλὰ εἶχον\FnGenericTemplate{ا}{\Parse{ماضی استمراری معلوم اخباری اول‌شخص مفرد} \Form{ἔχω} } ⸂γράψαι\FnGenericTemplate{ب}{\Parse{ماضی ساده معلوم مصدر} \Form{γράφω} } σοι⸃, ἀλλ’ οὐ θέλω διὰ μέλανος\FnGenericTemplate{پ}{\Parse{اضافی خنثی مفرد} \Form{μέλαν} \Gloss{جوهر} } καὶ καλάμου\FnGenericTemplate{ت}{\Parse{اضافی مذکر مفرد} \Form{κάλαμος} \Gloss{قلم} } σοι ⸀γράφειν·\FnGenericTemplate{ث}{\Parse{حال معلوم مصدر} \Form{γράφω} } 
\MainTextVerseMark{1}{14}ἐλπίζω δὲ εὐθέως ⸂σε ἰδεῖν⸃,\FnGenericTemplate{ا}{\Parse{ماضی ساده معلوم مصدر} \Form{ὁράω} } καὶ στόμα πρὸς στόμα λαλήσομεν.\FnGenericTemplate{ب}{\Parse{آینده معلوم اخباری اول‌شخص جمع} \Form{λαλέω} } 
\MainTextVerseMark{1}{15}Εἰρήνη σοι. ἀσπάζονταί σε οἱ φίλοι.\FnGenericTemplate{ا}{\Parse{فاعلی مذکر جمع} \Form{φίλος} \Gloss{دوست} } ἀσπάζου\FnGenericTemplate{ب}{\Parse{حال میانه امری دوم‌شخص مفرد} \Form{ἀσπάζομαι} } τοὺς φίλους\FnGenericTemplate{پ}{\Parse{مفعولی مذکر جمع} \Form{φίλος} \Gloss{دوست} } κατ’ ὄνομα. 
\end{document}